% Generated by analysis/scripts/export_ten_cases.py
\subsection*{统一结果卡}
\begingroup\small
\noindent\begin{tabularx}{\textwidth}{@{}p{24mm}X@{}}\toprule
技术/模式 & 3D vertical AND FeFET（2026）；2026 vertical AND二状态数字路径；16-cell写条带\\
逻辑粒度/聚合 & B\_S=128 Byte；B\_R=16 Byte。16 Byte完整权重，8个16-cell条带分两极性写验；1024事务覆盖矩阵\\
资源/相对R0 & 32横向读通路/4096感测节点；四层仅分装容量；16cell写驱动，128cell对照需8倍资源\\
服务口径 & 无周期维护；写后观察为工程预留，不是已证实必需等待\\
主导因素/选择 & $\pm$2 V/20 ns脉冲；100 ns写后观察预算已含最后返回，终态读验另计；观察占参考写时间46.2\%\\
关键对照 & 固定外围观察预算100/150 ns；128cell驱动资源对照；不由RAWD<100 ns推出零等待\\
\bottomrule\end{tabularx}
\vspace{3mm}
\begin{center}\begin{tabular}{@{}lrrr@{}}\toprule
成对条件情景 & $\rho$ (MB/s) & $\tau$ (MB/s) & $\mathrm{RI}^{*}$\\\midrule
短预算 & 22.71 & 11.27 & 2.016\\
参考（推荐） & 11.1 & 9.249 & 1.2\\
长预算 & 5.551 & 6.838 & 0.8118\\
\bottomrule\end{tabular}\end{center}
\noindent 参考原始占用：streaming 1.153e+04 ns；resident 1730 ns。
\par\noindent 范围为有限成对条件情景极值，不是物理不确定性包络。1 MB=$10^6$ Byte；整矩阵16384 Byte=16 KiB。源定位见本例证据笔记；公共基线shared\_baseline。
\endgroup
\clearpage
